We model a population of $M$ individuals making collective decisions on an NK
fitness landscape. Each individual~$i$ is represented by a binary string
$\mathbf{x}_i \in \{0,1\}^N$, partitioned into a \emph{shared portion}
$\mathbf{s} \in \{0,1\}^{N_s}$ (identical across all individuals, representing
laws) and an \emph{individual portion} $\mathbf{t}_i \in \{0,1\}^{N_t}$ (unique
and fixed, representing personal traits), where $N_s + N_t = N$.

\subsection{Fitness Landscape}

Fitness is computed according to the standard NK
model~\citep{kauffman1993origins}. Each bit position $j$ contributes a fitness
value $f_j$ that depends on the bit at position~$j$ and $K$ other positions
determined by a dependency matrix~$D$:
\begin{equation}
  F(\mathbf{x}_i) = \sum_{j=1}^{N} f_j\!\left(\mathbf{x}_i[j],\;
    \mathbf{x}_i[d_{j,1}],\ldots,\mathbf{x}_i[d_{j,K}]\right),
\end{equation}
where $d_{j,1},\ldots,d_{j,K}$ are the $K$ positions on which position~$j$
depends. Each $f_j$ maps a $(K\!+\!1)$-bit substring to a fitness value drawn
uniformly from $[-1, 1]$. Increasing $K$ increases the ruggedness of the
landscape, creating more local optima and making optimization harder.

\subsection{Cross-Dependency Parameter $\alpha$}

We control the coupling between shared and individual portions through a
cross-dependency fraction $\alpha \in [0,1]$. For each bit position~$j$, a
fraction~$\alpha$ of its $K$ dependencies are drawn from the \emph{other}
portion (shared bits depend on individual bits, and vice versa), while the
remaining $1-\alpha$ fraction are drawn from the same portion. When
$\alpha K$ is not an integer, we use stochastic rounding: the number of
cross-dependencies is $\lfloor \alpha K \rfloor + \text{Bernoulli}(\alpha K -
\lfloor \alpha K \rfloor)$.

At $\alpha=0$, the shared and individual portions are fitness-independent: laws
have the same effect regardless of individual traits. At $\alpha=1$, dependencies
are maximally entangled: the fitness effect of every bit depends on bits from
both portions. This captures the real-world phenomenon that the impact of
legislation depends on individual circumstances.

\subsection{Voting Process}

At each time step, the system selects $V$ bit positions uniformly at random from
the shared portion (we use $V=2$ throughout). This defines $2^V$ possible
proposals (all binary configurations of those positions). Each individual~$i$
evaluates all proposals by computing the fitness they would obtain under each
one, producing a utility matrix $U \in \mathbb{R}^{M \times 2^V}$.

A voting method $\mathcal{V}$ aggregates these preferences to select a winning
proposal, which is then applied to all individuals' shared portions. We study
eight methods:

\begin{itemize}
  \item \textbf{Plurality}: Each voter casts one vote for their highest-utility
    proposal; the proposal with the most votes wins.
  \item \textbf{Approval}: Each voter approves all proposals that improve their
    fitness over the status quo (or all tied-for-best if none improve); the most
    approved proposal wins.
  \item \textbf{Score}: Sum each voter's raw utility for each proposal; highest
    total wins.
  \item \textbf{Borda}: Each voter ranks proposals; points are assigned by rank
    and summed.
  \item \textbf{Instant-Runoff (IRV)}: Iteratively eliminate the proposal with
    the fewest first-place votes, redistributing those votes, until a majority
    winner emerges.
  \item \textbf{STAR}: Sum scores to find the top two proposals, then elect the
    one preferred head-to-head by more voters.
  \item \textbf{Minimax (Condorcet)}: Select the proposal whose worst pairwise
    margin of defeat is minimized.
  \item \textbf{Random Dictator}: A uniformly random voter's top choice is
    selected. Serves as a baseline.
\end{itemize}

\subsection{Representative Democracy}

In the representative variant, the population does not vote directly on
proposals. Instead, $C$ candidates are drawn from the population, each candidate
identifies their personally preferred proposal, and the entire population votes
(using the same voting methods above) on which \emph{candidate} to elect. The
winning candidate's preferred proposal is then implemented.

Candidates may be selected uniformly at random or via a softmax distribution over
current fitness with temperature $\tau$:
\begin{equation}
  P(\text{candidate } i) = \frac{\exp(F(\mathbf{x}_i)/\tau)}
    {\sum_{j=1}^{M}\exp(F(\mathbf{x}_j)/\tau)}.
\end{equation}
Low $\tau$ concentrates candidacy among high-fitness individuals (modeling
incumbency or elite advantage), while $\tau \to \infty$ recovers uniform
selection.
